\section{Comunicação, armazenamento e visualização}
\label{sec:comunicacao}

\subsection{Contrato de transmissão}

O envio é realizado pela ESP32 em rede Wi-Fi de 2,4~GHz, por
\texttt{POST /api/readings}. A API valida tipos, limites e coerência
dos indicadores de falha; um registro aceito e persistido retorna HTTP~201.
Payloads incompatíveis retornam 422. Falha de gravação do CSV retorna
503, sem substituir a última leitura aceita.
O Código~\ref{cod:payload} utiliza os nomes efetivamente exigidos
pelo contrato atual.

\begin{lstlisting}[caption={Payload ilustrativo compatível com o contrato da API},label={cod:payload}]
{
  "device_id": "ms-monitor-g3",
  "uptime_ms": 125000,
  "window_ms": 2500,
  "temperature_c": 24.8,
  "humidity_pct": 59.2,
  "dht_ok": true,
  "probe_temperature_c": 25.25,
  "probe_ok": true,
  "vibration_do": 0,
  "vibration_transitions": 12,
  "embedded_diagnosis": {
    "source": "esp32",
    "model_version": "embedded-rf-6b0d1696547d",
    "risk_level": "low",
    "risk_score": 0.12,
    "confidence": 0.8,
    "inference_us": 123
  }
}
\end{lstlisting}

Todos os valores desse exemplo são \textbf{ilustrativos}. Em particular,
123~$\mu$s não é uma duração medida na placa. O objeto exemplifica o
contrato, e não um ensaio físico. A diferença térmica é calculada
internamente para a IA e não constitui um campo próprio no POST atual.

O campo \texttt{uptime\_ms} representa tempo desde a inicialização da placa,
não horário civil. O servidor acrescenta \texttt{received\_at} em UTC.
Esse carimbo corresponde ao recebimento, sujeito ao atraso de comunicação,
e não ao instante exato de aquisição. A API não executa outra inferência:
apenas valida e prepara o diagnóstico recebido para exibição.

\subsection{Consulta e persistência}

\begin{table}[htbp]
\centering
\caption{Interfaces HTTP implementadas}
\label{tab:api}
\small
\begin{tabularx}{\textwidth}{L{5.0cm}Y}
\toprule
Rota & Finalidade\\
\midrule
\texttt{GET /health} & Saúde da API, independente do estado do dispositivo.\\
\texttt{POST /api/readings} & Recebimento e persistência de uma leitura e de seu diagnóstico embarcado.\\
\texttt{GET /api/readings/latest} & Última leitura, idade, estado de atualização e diagnóstico recebido.\\
\texttt{GET /api/readings?limit=300} & Histórico, em ordem cronológica, com limite de 1 a 3600 registros.\\
\texttt{GET /api/readings/export} & Download do CSV completo até o instante da solicitação.\\
\bottomrule
\end{tabularx}
\source{elaboração própria a partir de \texttt{backend/app.py}}
\end{table}

O arquivo padrão é \path{backend/data/readings.csv}. Cada POST aceito
gera uma linha com as medições, os indicadores de validade e os campos
\path{ai_source}, \path{ai_model_version}, \path{ai_risk_level},
\path{ai_risk_score}, \path{ai_confidence} e
\path{ai_inference_us}. Campos ausentes ficam em branco no CSV.
As colunas de IA registram previsões, e não rótulos independentes de
contaminação para treinamento supervisionado.

O histórico em memória é limitado aos últimos 3600 registros e é
recuperado do CSV ao iniciar a API. O arquivo completo continua crescendo,
portanto a coleta exige acompanhamento do espaço em disco e cópias de
segurança. Deve haver uma única instância da API, com um único worker,
por arquivo. Após reiniciar, o histórico permanece disponível, mas o
estado da última leitura aguarda um novo envio.

O envelope de consulta diferencia \texttt{waiting}, antes de um POST
aceito pela instância, \texttt{online}, para leitura com idade inferior
a 10~s, e \texttt{stale}, a partir desse intervalo. O frescor é calculado
por tempo monotônico no servidor. Um dispositivo com dados recentes
ainda pode reportar falha de um sensor; por isso, conectividade e
validade das medições são informações distintas.

\subsection{Apresentação e limites de comunicação}

O dashboard exibe temperatura do ar, temperatura da sonda, UR,
transições por janela, última atualização e resultado embarcado com
versão e tempo de inferência. A diferença térmica é uma entrada
interna do modelo, não um card dedicado na interface examinada.
O card de ventilação permanece como \textit{Não avaliado}.

Para os valores ilustrativos do Código~\ref{cod:payload}, a interface
apresentaria 24,8\celsius\ de temperatura do ar, 59,2\% de UR,
25,25\celsius\ na sonda e 12 transições em uma janela de 2,5~s,
com DO igual a zero. A classe \texttt{low} apareceria como
\textit{Baixo risco}, com índice de 12\%, confiança de 80\% e
inferência de 123~$\mu$s. Essa correspondência exemplifica a
formatação dos dados recebidos, sem transformar o exemplo em
resultado de ensaio ou o índice em confirmação de contaminação.

A última leitura é consultada aproximadamente a cada segundo; o
histórico é sincronizado de forma independente. O modo \textit{Ao vivo}
prioriza a coleta atual, enquanto o histórico e o CSV preservam
registros anteriores. Uma pausa de recepção não é preenchida com
valores artificiais. Na indisponibilidade da API, os valores anteriores
continuam visíveis com indicação de desatualização.

A API local usa HTTP e não implementa autenticação de dispositivos.
CORS restringe origens de navegador, mas não autentica POSTs enviados
pela ESP32. A arquitetura não deve ser apresentada como solução de
segurança para acesso público à internet. Credenciais de Wi-Fi ficam
em um \texttt{config.h} local e não integram o repositório ou este relatório.
\FloatBarrier
